\section{Market, contract and data}
\label{sec:market}

\subsection{The Turkish wholesale electricity market}

Wholesale electricity in Türkiye clears through EPİAŞ, the market operator. The
day-ahead market produces an hourly clearing price, the \emph{Piyasa Takas
Fiyatı} (PTF), for every delivery hour of the following day \citep{EPIAS2025}.
PTF is the reference price for physical settlement and the underlying taken here
for option valuation. Prices are quoted in Turkish lira per megawatt-hour and are
subject to administratively set ceiling and floor levels revised periodically
\citep{EPDK2024}.

The organised forward segment, VEP, lists monthly baseload futures. A contract
for delivery month $m$, denoted \texttt{EBM}$mm$\texttt{yy}, settles against the
arithmetic mean of PTF over every delivery hour of that month. This averaging
structure is the first source of the identification problem addressed here. A
monthly baseload quote constrains the \emph{average} of $720$ to $744$ hourly
forward prices and says nothing about how that average is distributed across
hours. Two curves with identical monthly means but different intraday shapes are
equally consistent with the same quote, yet imply materially different values for
an option written on a single delivery hour.

Two features make the market a useful laboratory for these questions. Renewable
penetration has grown quickly, so residual demand fluctuates over a wide range
and the merit-order mechanism transmitting it into price is strong
\citep{SirinYilmaz2020}. And the lira has depreciated substantially over the
period from which historical parameters must be drawn, so nominal price series
mix genuine dynamics with monetary trend. Consequences of the second feature are
taken up in Section~\ref{sec:limitations}.

\subsection{The contract}

The instrument valued here is a European call on the PTF observed at a single
expiry hour $T$,
\begin{equation}
  V_0 \;=\; e^{-rT}\,\E^{\Q}\!\left[(P_T-K)^+\right],
  \label{eq:payoff}
\end{equation}
where $P_T$ is the hourly clearing price at $T$, $K$ the strike and $r$ a flat
continuously compounded annual rate. This is not the monthly baseload contract
that VEP quotes. The distinction is enforced in the implementation: monthly
quotes are represented by an object that raises an error if accessed as a point
forward, since conflating the two is the shortcut the data structure most invites.

\subsection{Valuation date and samples}

The valuation timestamp is 2025-12-31 20:00 UTC, the final hour of the historical
sample, at which the observed spot is $2{,}917.78$ \TRY. Six monthly baseload
quotes were observed at that date, for delivery months February through July 2026
(Table~\ref{tab:quotes}). January 2026 carries no quote. This is not an avoidable
data gap: the nearest delivery month is the one whose option values matter most
for short-dated contracts, and it is precisely the month the strip does not
price. Section~\ref{sec:anchor} treats it explicitly.

Two further samples enter. The regime covariate is a standardised residual demand
series $z_t$ at hourly frequency, standardised using the training scaler of the
underlying estimation exercise ($N=61{,}361$). The realised evaluation uses hourly
PTF for January through July 2026, $5{,}088$ hours, retrieved from the EPİAŞ
transparency platform after the fact. The evaluation sample is strictly posterior
to the valuation timestamp; no observation from it enters calibration, parameter
estimation or model selection. This is verified structurally rather than by
assertion, through a guard test that fails if any fitting routine is exposed to a
timestamp at or beyond the valuation cutoff.

\begin{table}[htbp]
\centering
\caption{Observed VEP monthly baseload quotes at the valuation date
(2025-12-31), in \TRY.}
\label{tab:quotes}
\small
\begin{tabular}{@{}llrr@{}}
\toprule
Contract & Delivery month & Delivery hours & Quote \\
\midrule
EBM0226 & February 2026 & 672 & 2{,}900.99 \\
EBM0326 & March 2026    & 744 & 2{,}556.31 \\
EBM0426 & April 2026    & 720 & 2{,}500.66 \\
EBM0526 & May 2026      & 744 & 2{,}506.25 \\
EBM0626 & June 2026     & 720 & 2{,}245.33 \\
EBM0726 & July 2026     & 744 & 3{,}558.19 \\
\midrule
\multicolumn{2}{@{}l}{January 2026} & 744 & \emph{unquoted} \\
\bottomrule
\end{tabular}
\end{table}

The strip carries information worth noting before any modelling. The market
priced a pronounced decline from February to June followed by a sharp July
rebound of roughly $58\%$ over the June level, a shape consistent with
cooling-driven summer peaks and high spring hydro and solar availability. The
hourly curve must therefore accommodate a large discontinuity at the June--July
boundary originating in the data rather than in the interpolation method, a point
verified in Section~\ref{sec:curvefit}.
